% !PRINTED 307
\weberpage{307}

%<<CONT
$\varrho$ durchläuft in (7) und (8) entweder die Reihe der quadratischen Reste oder die der Nichtreste, so daß man für $p = 7$ oder 11 \gesperrt{je zwei Gruppen vom Index $p$ erhält.}

Daß die in (3), (7), (8) zusammengestellten Substitutionssysteme wirklich Gruppen konstituieren, läßt sich durch direkte Zusammensetzung auf verschiedene Arten nachweisen. Einfacher gelangt man zu diesem Beweise aber dadurch, daß man Funktionen der $v_z$, § 78, (7), bildet, die durch die Substitutionen dieser Systeme ungeändert bleiben, und durch die Substitutionen von $\mathfrak{L}_0$ überhaupt nur $p$ verschiedene Werte erhalten.

Die Systeme (3), (7), (8) lassen sich, wie aus ihrer Entstehungsweise hervorgeht, durch wiederholte Anwendung von $B, C, U, U'$ zusammensetzen, wenn $U, U'$ irgend zwei spezielle Substitutionen $W, W'$ sind. Für $p = 5$ gehört $U^2$ und für $p = 7, 11$ gehört $UB$ zu den $W'$, so daß $U'$ noch weggelassen werden kann. Wählen wir $U$ irgendwie beliebig, und nehmen für die primitive Wurzel $g$ in $C$ für $p = 5, 7, 11$ bzw. $2, -2, 2$, so können wir die Gruppe (3) zusammensetzen aus:
\[
B = \begin{pmatrix} 0, & 1 \\ -1, & 0 \end{pmatrix}, \quad C = \begin{pmatrix} 2, & 0 \\ 0, & 3 \end{pmatrix}, \quad U = \begin{pmatrix} 1, & 1 \\ 2, & 3 \end{pmatrix} \quad (p = 5),
\]
die Gruppe (7) aus:
\[
B = \begin{pmatrix} 0, & 1 \\ -1, & 0 \end{pmatrix}, \quad C = \begin{pmatrix} 2, & 0 \\ 0, & 4 \end{pmatrix}, \quad U = \begin{pmatrix} 1, & \pm 1 \\ \pm 2, & 3 \end{pmatrix} \quad (p = 7),
\]
die Gruppe (8) aus:
\[
B = \begin{pmatrix} 0, & 1 \\ -1, & 0 \end{pmatrix}, \quad C = \begin{pmatrix} 2, & 0 \\ 0, & 6 \end{pmatrix}, \quad U = \begin{pmatrix} 1, & \pm 1 \\ \pm 1, & 2 \end{pmatrix} \quad (p = 11),
\]
wo in den beiden letzten Fällen aus den doppelten Vorzeichen die oben erwähnten zwei verschiedenen Gruppen entspringen.

Wir stellen nun die diesen Substitutionen entsprechenden Vertauschungen der Indizes $z$ zusammen [§ 79, (3)].

$p = 5$:
\[
\begin{array}{ccccccl}
z = \infty, & 0, & 1, & 2, & 3, & 4 & \\
0, & \infty, & 4, & 2, & 3, & 1 & (B)\\
\infty, & 0, & 4, & 3, & 2, & 1 & (C)\\
4, & 1, & 2, & 0, & \infty, & 3 & (U).
\end{array}
\]
Es werden also durch $B, C, U$ die Indexpaare
\[
(\infty, 0), \quad (1, 4), \quad (2, 3)
\]
%CONT>>
